% !TeX encoding = UTF-8
% 吉田第六版：中译本第 0 章，四节正文与参考文献整理。
\SetChapterInfo{subtitle={集合、拓扑、测度与线性结构},keywords={选择原理; 紧性; 函数逼近; 积分; 商空间},prereq={集合与映射、微积分、线性代数},goals={理解选择原理的等价性; 掌握紧性与类别方法; 证明积分收敛与交换定理; 建立线性算子和商空间}}
\chapter{预备知识}\label{ch:yosida-00}
\ChapterSource{本章保留吉田耕作第六版第零章的四节结构，原书位置为中译本与英文版印刷页 1--22，参见 \cite{yosida-zh,yosida-en}。正文按原书主题重组为讲义，并补充选择原理的等价性证明、积分构造及代数应用；使用现代测度术语，涉及原书与本稿的记号差别时另行说明。本章参考作者及其总书目条目已核对，正文不是逐句翻译。}

\section{集合论}\label{sec:yosida-00-01}
\noindent\textit{原书位置：中译本与英文版本节均从印刷第 1 页起。}
% !TeX encoding = UTF-8
泛函分析中常须把一个局部构造延拓到整个空间。例如，线性无关集能否扩充为基，或者已定义的线性泛函能否继续延拓？这类问题常通过“按包含关系排列所有候选对象，再取一个极大对象”解决。本节建立这一方法的集合论基础。

以下在 Zermelo--Fraenkel 集合论（ZF）中讨论；选择公理作为额外命题单独提出。所有偏序集、索引族及线性空间的底集都是集合。记 $\mathcal P(X)$ 为 $X$ 的幂集，$\varnothing$ 为空集。三个选择原理之间的证明不能预先借用其中另一个；为此，先交代良序与超限递归所需的基础工具。

\subsection{偏序、链与极大元}
\begin{definition}[偏序与全序]\label{def:y0-partial-order}
集合 $P$ 上的二元关系 $\le$ 称为\emph{偏序}，若它满足：
\begin{enumerate}
  \item 自反性：$x\le x$，对每个 $x\in P$ 成立；
  \item 反对称性：$x\le y$ 且 $y\le x$ 蕴含 $x=y$；
  \item 传递性：$x\le y$ 且 $y\le z$ 蕴含 $x\le z$。
\end{enumerate}
称 $(P,\le)$ 为偏序集，记 $x<y$ 表示 $x\le y$ 且 $x\ne y$。若任意两个元素都可比较，即 $x\le y$ 或 $y\le x$，则此偏序称为\emph{全序}。
\end{definition}

\begin{definition}[链、上界、最大元与极大元]\label{def:y0-chain}
设 $(P,\le)$ 为偏序集，$A\subseteq P$。
\begin{enumerate}
  \item 若 $C\subseteq P$ 中任意两元素可比较，则称 $C$ 为\emph{链}；允许 $C=\varnothing$。
  \item 若 $u\in P$ 且每个 $a\in A$ 都满足 $a\le u$，则称 $u$ 为 $A$ 在 $P$ 中的\emph{上界}。上界不必属于 $A$。上界中的最小元称为 $A$ 的\emph{上确界}，记为 $\sup A$。
  \item 若 $a_*\in A$ 且每个 $a\in A$ 都有 $a\le a_*$，则称 $a_*$ 为 $A$ 的\emph{最大元}。
  \item 若 $m\in P$ 且 $m\le x$ 蕴含 $x=m$，则称 $m$ 为 $P$ 的\emph{极大元}。
\end{enumerate}
反向不等式分别定义下界、下确界、最小元与极小元。讨论子集的极大元时，使用从 $P$ 限制得到的偏序。
\end{definition}

最大元必须支配所有元素，而极大元只要求不存在严格更大的元素。最大元若存在则唯一，并且是极大元；极大元可以不唯一。在全序集中，极大元也是最大元。空链的每个上界条件都是空条件，故它的上界恰为 $P$ 的全部元素。

\begin{example}[包含序与整除序]
幂集 $\mathcal P(X)$ 在包含关系 $\subseteq$ 下是偏序集。若 $X$ 至少有两个元素，它通常不是全序集。任意子族 $\mathcal A\subseteq\mathcal P(X)$ 的上确界为 $\bigcup\mathcal A$，下确界为 $\bigcap\mathcal A$；这里约定空族的交为 $X$。

在 $P=\{\{1\},\{2\}\}$ 的包含序下，两个元素均极大，却没有最大元。正整数集在整除关系下也是偏序集；例如 $2$ 与 $3$ 不可比较。通常的数值大小和整除关系是两种不同的偏序。
\end{example}

\subsection{良序与超限递归的基础}
\begin{definition}[良序与初始段]\label{def:y0-well-order}
全序集 $(X,\le)$ 称为\emph{良序集}，若每个非空子集 $A\subseteq X$ 都有最小元。对 $x\in X$，集合
\[
 X_{<x}=\{y\in X:y<x\}
\]
称为由 $x$ 决定的严格初始段。一般地，$I\subseteq X$ 若满足 $x\in I$ 且 $y<x$ 时 $y\in I$，则称 $I$ 为初始段。
\end{definition}

自然数集在通常次序下良序。整数集不是良序，因为整个整数集没有最小元；实数集不是良序，例如 $(0,1)$ 没有最小元。有限全序集总是良序。\emph{良序原则断言可给一个集合另赋某种良序，并非断言它原有的次序就是良序。}

\begin{proposition}[良序归纳]\label{prop:y0-well-order-induction}
设 $X$ 良序。若性质 $Q(x)$ 满足：对任意 $x\in X$，所有 $y<x$ 满足 $Q(y)$ 时，$Q(x)$ 也成立，则每个 $x\in X$ 都满足 $Q(x)$。
\end{proposition}
\begin{proof}
若反例集非空，取其中的最小元 $x_0$。所有 $y<x_0$ 都不是反例，由假设 $Q(x_0)$ 成立，矛盾。
\end{proof}

为精确说明“逐个选择，直到选完”何时停止，使用序数记号。\emph{序数}是由成员关系 $\in$ 严格良序的传递集合；传递指 $u\in\alpha$ 蕴含 $u\subseteq\alpha$。序数的严格次序记为 $\beta<\alpha$，即 $\beta\in\alpha$。最小序数为 $0=\varnothing$，后继序数为 $\alpha+1=\alpha\cup\{\alpha\}$；非零且不是后继的序数称为极限序数。自然数的序数集合记为 $\omega$。一族序数的并仍为序数，且是该族的上确界。

每个良序集都与唯一序数保序同构，此序数称为其\emph{序型}。这里“保序同构”指保持并反映次序的双射。其构造是沿良序定义
\[
 r(x)=\{r(y):y<x\}.
\]
良序归纳表明各 $r(x)$ 是序数，且 $y<x$ 当且仅当 $r(y)\in r(x)$；像集是传递的序数。唯一性也由良序归纳得到：两个保序的编号在每个位置都有相同的前驱编号，因而相等。这一构造使用下面的递归定理，它在 ZF 内成立，不需要选择公理。

\begin{lemma}[沿良序的递归]\label{lem:y0-transfinite-recursion}
设 $X$ 为良序集。若一条确定规则对每个 $x\in X$ 及已定义在 $X_{<x}$ 上的函数 $g$ 指定唯一集合值 $F(x,g)$，则存在唯一函数 $f$，定义域为 $X$，且
\[
 f(x)=F\bigl(x,f|_{X_{<x}}\bigr)\qquad(x\in X).
\]
在序数上使用此定理称为超限递归；在极限阶段，已有函数取并，得到此前全部取值。
\end{lemma}
\begin{proof}
称定义在某个初始段上并满足递归式的函数为近似解。任意两个近似解在公共定义域上一致：若有分歧，取最早分歧点；该点之前的函数相同，规则的唯一性迫使此点取值也相同，矛盾。

用良序归纳证明每个 $X_{<x}$ 上有唯一近似解。若所有更早的初始段都有解，对每个 $y<x$ 将 $X_{<y}$ 上的解延拓一步，定义 $y$ 处的取值为 $F(y,g)$。由替代公理，这些唯一确定的延拓组成一个集合；由相容性，其并给出 $X_{<x}$ 上的解。空初始段的解为空函数。最后，对所有 $x\in X$ 再将这些解延拓一步并取并，得到整个 $X$ 上的 $f$。相容性同时证明唯一性。这里收集的是唯一确定的对象，所用的是替代公理，而非选择公理。
\end{proof}

\begin{lemma}[Hartogs 引理]\label{lem:y0-hartogs}
对任意集合 $X$，存在序数 $\theta$，使得不存在从 $\theta$ 到 $X$ 的单射。此结论可在 ZF 中证明。
\end{lemma}
\begin{proof}
考虑 $X$ 的各个子集上的所有良序关系。每一项可表示为一对 $(A,R)$，其中 $A\subseteq X$，$R\subseteq X\times X$，故这些对象构成一个集合。对每个这样的良序取唯一序型，由替代公理得到一个序数集合 $T$。令
\[
 \theta=\bigcup_{\alpha\in T}(\alpha+1).
\]
这是序数，且对所有 $\alpha\in T$ 都有 $\alpha<\theta$。若有单射 $j:\theta\to X$，把 $\theta$ 的良序搬到像集 $j[\theta]$ 上，得到 $X$ 的一个子集上的良序，其序型正是 $\theta$。于是 $\theta\in T$，从而 $\theta<\theta$，矛盾。
\end{proof}

Hartogs 引理给出递归选择不能无限延长的集合论界限。只说“每次选一个，最后必能选完”并不构成证明；若只按自然数编号，更无法遍历任意不可数集合。

\subsection{选择公理、良序原则与 Zorn 引理}
\begin{definition}[选择函数与选择公理]\label{def:y0-choice}
设 $(A_i)_{i\in I}$ 为一族非空集合。函数 $f:I\to\bigcup_{i\in I}A_i$ 若满足 $f(i)\in A_i$，称为这族集合的\emph{选择函数}。\emph{选择公理}（AC）断言：每一族非空集合都有选择函数。等价地，每一这样的笛卡尔积都非空：
\[
 \prod_{i\in I}A_i
 =\left\{f:I\to\bigcup_{i\in I}A_i:
          f(i)\in A_i\text{ 对所有 }i\in I\right\}\ne\varnothing.
\]
若 $I=\varnothing$，空函数即为选择函数，此时乘积是单元素集 $\{\varnothing\}$。
\end{definition}

有限族非空集合的选择函数可用有限归纳在 ZF 中构造；一般无限族的结论是额外公理。若每个 $A_i$ 的元素由一条统一规则唯一确定，例如非空自然数子集的最小元，也不需要选择公理。选择公理保证函数存在，不提供可计算的选择算法。

\begin{theorem}[良序原则]\label{thm:y0-well-order-principle}
每个集合 $X$ 都可以赋予一个良序关系。空集的空关系也视为良序。
\end{theorem}

\begin{theorem}[Zorn 引理]\label{thm:y0-zorn}
设 $(P,\le)$ 是非空偏序集。如果每条链 $C\subseteq P$ 都在 $P$ 中有上界，则 $P$ 至少有一个极大元。
\end{theorem}

Zorn 引理只要求上界存在，不要求上界唯一，也不要求上确界存在。若把空链也包含在条件中，“每条链都有上界”本身已蕴含 $P\ne\varnothing$；若只考察非空链，则必须另列 $P\ne\varnothing$。

\begin{theorem}[三种选择原理的等价性]\label{thm:y0-choice-equivalence}
在 ZF 集合论中，以下三命题等价：选择公理、良序原则、Zorn 引理。
\end{theorem}
\begin{proof}
依次证明
\[
 \text{选择公理}\ \Longrightarrow\ \text{良序原则}
 \ \Longrightarrow\ \text{Zorn 引理}
 \ \Longrightarrow\ \text{选择公理}.
\]

\emph{第一步：选择公理推出良序原则。}
若 $X=\varnothing$，结论直接成立。设 $X\ne\varnothing$，对所有非空子集应用选择公理，固定函数
\[
 c:\mathcal P(X)\setminus\{\varnothing\}\longrightarrow X,
 \qquad c(A)\in A.
\]
由 Hartogs 引理取不能单射到 $X$ 的序数 $\theta$。取一个不在 $X$ 中的标记 $\ast$，沿 $\theta$ 递归定义
\[
 R_\alpha=X\setminus
       \{x_\beta:\beta<\alpha,\ x_\beta\in X\},\qquad
 x_\alpha=
 \begin{cases}
   c(R_\alpha),&R_\alpha\ne\varnothing,\\
   \ast,&R_\alpha=\varnothing.
 \end{cases}
\]
若每个 $\alpha<\theta$ 都有 $R_\alpha\ne\varnothing$，则 $\alpha\mapsto x_\alpha$ 是 $\theta$ 到 $X$ 的单射，因为每次所取元素均未在此前出现。这与 $\theta$ 的定义矛盾。

故存在最小的 $\gamma<\theta$ 使 $R_\gamma=\varnothing$。在 $\gamma$ 之前，取值两两不同，且已遍历 $X$，于是 $\alpha\mapsto x_\alpha$ 是 $\gamma$ 到 $X$ 的双射。定义 $x_\alpha\prec x_\beta$ 当且仅当 $\alpha<\beta$，便把序数 $\gamma$ 的良序搬到了 $X$ 上。

\emph{第二步：良序原则推出 Zorn 引理。}
设 $P$ 非空且每条链都有上界。由良序原则，将 $P$ 的全部元素无重复地写作 $(p_\alpha)_{\alpha<\kappa}$，其中 $\kappa$ 是某个非零序数。按这个编号逐个检查元素，只在新元素严格大于所有已保留元素时将它保留。精确地，令 $C_0=\varnothing$，并递归定义
\[
 C_{\alpha+1}=
 \begin{cases}
 C_\alpha\cup\{p_\alpha\},
       &x<p_\alpha\text{ 对每个 }x\in C_\alpha,\\
 C_\alpha,&\text{否则},
 \end{cases}
 \qquad
 C_\lambda=\bigcup_{\alpha<\lambda}C_\alpha
 \quad(\lambda\text{ 为极限序数}).
\]
良序归纳表明各 $C_\alpha$ 是链，且随 $\alpha$ 递增；极限阶段的并仍为链，因为任意两个元素已共同处于某个更早阶段。于是 $C_\kappa$ 是链，并且 $p_0$ 在第一步被保留，所以它非空。

取 $C_\kappa$ 的一个上界 $u\in P$。若 $u$ 不是极大元，则存在 $v\in P$ 使 $u<v$。设 $v=p_\beta$。在检查 $p_\beta$ 的阶段，对所有 $x\in C_\beta$ 都有 $x\le u<v$，所以 $v$ 被保留，属于 $C_\kappa$。上界条件又给出 $v\le u$，与 $u<v$ 矛盾。因此 $u$ 为极大元。

\emph{第三步：Zorn 引理推出选择公理。}
设 $(A_i)_{i\in I}$ 为非空集合族，令 $U=\bigcup_{i\in I}A_i$。考虑全部部分选择函数：
\[
 P=\{f:J\to U:J\subseteq I,\ f(i)\in A_i\ (i\in J)\},
\]
把函数看作图，即 $I\times U$ 的子集，以图的包含关系作为偏序。因空函数属于 $P$，此偏序集非空。

若 $\mathcal C\subseteq P$ 是链，令 $g=\bigcup_{f\in\mathcal C}f$。这是函数：同一指标若在两个成员的定义域中出现，这两个成员可比较，较大的函数延拓较小的函数，取值必然一致。且 $g(i)\in A_i$，所以 $g\in P$，并为 $\mathcal C$ 的上界。空链的并为空函数，也满足要求。

由 Zorn 引理存在极大部分选择函数 $f_*$。若其定义域 $J_*$ 不等于 $I$，取 $i_0\in I\setminus J_*$，再取一个 $a_0\in A_{i_0}$，则
\[
 f_*\cup\{(i_0,a_0)\}
\]
是严格延拓 $f_*$ 的部分选择函数，与极大性矛盾。因此 $J_*=I$，$f_*$ 即为所求选择函数。这里仅从一个非空集合取一个元素，不是在所有 $A_i$ 中同时作选择，故没有循环使用选择公理。
\end{proof}

上述闭环已经证明任意两命题互相蕴含。下面再给出两个常用的直接证明，帮助识别“良序化”和“取极大延拓”这两种构造。

\begin{proposition}[良序原则直接推出选择公理]\label{prop:y0-wo-ac}
若每个集合都可良序化，则每族非空集合都有选择函数。
\end{proposition}
\begin{proof}
将 $U=\bigcup_{i\in I}A_i$ 良序化，定义 $f(i)$ 为 $A_i$ 在这一良序下的最小元。非空性保证最小元存在，良序性保证唯一性；替代公理将这些唯一取值收集为函数。若 $I$ 为空，取空函数。
\end{proof}

\begin{proposition}[Zorn 引理直接推出良序原则]\label{prop:y0-zorn-wo}
若 Zorn 引理成立，则每个集合 $X$ 都可良序化。
\end{proposition}
\begin{proof}
令 $P$ 为 $X$ 的各个子集上的所有良序对 $(A,\le_A)$。规定
\[
 (A,\le_A)\preccurlyeq(B,\le_B)
\]
当且仅当 $A\subseteq B$，$\le_B$ 在 $A$ 上限制为 $\le_A$，且 $A$ 是 $B$ 的初始段。这是偏序，空集上的良序对使 $P$ 非空。

对非空链 $\mathcal C$，取底集的并 $B$ 和各个序关系的并 $\le_B$。任意有限个元素都共同处于某个链成员中，故 $\le_B$ 为全序。它还是良序：设 $S\subseteq B$ 非空，取 $s\in S$ 及含 $s$ 的链成员 $A$，令 $a$ 为 $S\cap A$ 在 $A$ 中的最小元。若 $t\in S\setminus A$，含 $t$ 的链成员必延拓 $A$，且 $A$ 是它的初始段，所以 $a<t$。于是 $a$ 是 $S$ 的最小元。每个链成员均为 $B$ 的初始段，故 $(B,\le_B)$ 为链的上界。空链由空良序对作为上界。

由 Zorn 引理取极大良序对 $(A_*,\le_*)$。若 $A_*\ne X$，取 $x\in X\setminus A_*$，把 $x$ 放在 $A_*$ 的全部元素之后。新次序仍是良序，且 $A_*$ 是其初始段，从而得到严格延拓，矛盾。所以 $A_*=X$。
\end{proof}

\begin{remark}[为什么需要初始段条件]
上一个证明若只要求两个良序的限制相容，其链的并未必良序。例如 $A_n=\{-n,\ldots,-1,0\}$ 各自在通常大小次序下良序，次序逐步延拓，但它们的并没有最小元。初始段条件禁止在已有元素之前不断插入新元素，正是并集良序性的关键。
\end{remark}

\subsection{极大链原理与泛函分析中的应用}
\begin{theorem}[Hausdorff 极大链原理]\label{thm:y0-maximal-chain}
每个偏序集的任意一条链都包含于一条按包含关系极大的链。此原理在 ZF 中与 Zorn 引理等价。
\end{theorem}
\begin{proof}
先假设 Zorn 引理。固定链 $C_0\subseteq P$，考虑全部包含 $C_0$ 的链组成的集合 $\mathcal Q$，以包含关系排序。它非空。$\mathcal Q$ 中一条链的并仍为 $P$ 中的链：任意两元素分别来自两个可比较的成员，故共同属于较大者。这个并包含 $C_0$，所以是 $\mathcal Q$ 中的上界；空链可用 $C_0$ 作上界。Zorn 引理给出所求极大链。

反过来，设 $P$ 非空且每条链有上界。用极大链原理将空链扩充为极大链 $C$，取其上界 $u$。由于 $C\cup\{u\}$ 仍为链，极大性迫使 $u\in C$。若 $u<v$，则 $C\cup\{v\}$ 也为链，故 $v\in C$；但 $u$ 是 $C$ 的上界，得到 $v\le u$，矛盾。所以 $u$ 为极大元。
\end{proof}

\begin{theorem}[线性无关集扩充为 Hamel 基]\label{thm:y0-hamel-basis}
假设选择公理，设 $V$ 是域 $\mathbb K$ 上的线性空间，$L\subseteq V$ 为线性无关集。则存在包含 $L$ 的 Hamel 基 $B$；也就是说，每个向量都唯一表示为 $B$ 中有限个向量的线性组合。特别地，每个线性空间都有 Hamel 基。
\end{theorem}
\begin{proof}
令 $P$ 为 $V$ 中全部包含 $L$ 的线性无关子集，以包含关系为偏序。$L\in P$，故 $P$ 非空。对非空链 $\mathcal C\subseteq P$，令 $M=\bigcup_{A\in\mathcal C}A$。任何有限个 $M$ 中的向量都共同处于某个链成员中：对有限个包含它们的成员逐次取较大者即可。因此它们的线性关系只能是平凡关系，$M$ 线性无关，并包含 $L$，故 $M\in P$。空链以 $L$ 为上界。

由 Zorn 引理取极大元素 $B$。若 $v\notin\operatorname{span}B$，则 $B\cup\{v\}$ 仍线性无关：任何涉及 $v$ 的有限线性关系，若 $v$ 的系数非零便把 $v$ 表示为 $B$ 的有限线性组合；若系数为零则由 $B$ 的线性无关性知其余系数全零。这样得到严格更大的成员，与极大性矛盾。因此 $\operatorname{span}B=V$。线性无关性保证表达唯一。取 $L=\varnothing$ 即得基的存在性；零空间的基为空集。
\end{proof}

\begin{remark}[代数基与拓扑基]
Hamel 基只涉及有限线性组合，不涉及收敛。Banach 空间中的 Schauder 基要求向量有收敛的级数展开；Hilbert 空间中的完全正交系也涉及内积与闭线性张成。Hamel 基存在定理不保证任意 Banach 空间具有 Schauder 基。
\end{remark}

对后续的 Hahn--Banach 延拓，候选对象则是定义在不同子空间上的线性泛函，以“定义域扩大且取值相容”排序。链的并给出一致的线性泛函，极大性再配合单步延拓将定义域推到整个空间。因此使用 Zorn 引理时，应依次明确：候选对象形成集合，偏序及非空性成立，每条链有合法上界，极大对象具有目标性质。

\subsection{补充习题与解答}
\begin{exercise}[极大元可以不唯一]
设 $X$ 为至少有两个元素的集合，$P=\mathcal P(X)\setminus\{X\}$，以包含关系排序。求 $P$ 的全部极大元，并判断 $P$ 是否有最大元。
\end{exercise}
\begin{solution}
极大元恰为 $X\setminus\{x\}$（$x\in X$）。这种集合再加入任何缺少的元素就变成 $X$，故极大。若 $A\in P$ 缺少两个不同元素 $x,y$，则 $A\cup\{x\}$ 仍是真子集，故 $A$ 不极大。若有最大元，它必须包含每个单元素子集，从而等于 $X$，与属于 $P$ 矛盾。
\end{solution}

\begin{exercise}[有界的数列并不足够]
在 Zorn 引理中，能否将“每条链都有上界”改为“每个可数链都有上界”？提示：考虑全体可数序数 $\omega_1$ 在通常序数次序下的偏序，并在选择公理下讨论。
\end{exercise}
\begin{solution}
不能。在选择公理下，可数个可数集合的并可数，所以任意可数子集 $C\subseteq\omega_1$ 的序数上确界 $\bigcup C$ 仍是可数序数，属于 $\omega_1$，并是其上界；空集的并为 $0$。但任意 $\alpha<\omega_1$ 的后继 $\alpha+1$ 仍在 $\omega_1$ 中，故没有极大元。这里可数并的结论使用了选择公理所蕴含的可数选择，不能在讨论 ZF 中的等价性时不加说明地借用。
\end{solution}

\begin{exercise}[已指定的局部选择可以延拓]
假设 Zorn 引理。对非空集合族 $(A_i)_{i\in I}$，若已给定 $J_0\subseteq I$ 上的部分选择函数 $f_0$，证明它能延拓为整个 $I$ 上的选择函数。
\end{exercise}
\begin{solution}
将定理 \ref{thm:y0-choice-equivalence} 第三步中的偏序集换为所有延拓 $f_0$ 的部分选择函数。它包含 $f_0$，因而非空；非空链的并仍延拓 $f_0$，空链以 $f_0$ 作上界。由 Zorn 引理取极大延拓，若有未定义指标，选其中一个指标及其集合中的一个元素即可继续延拓，矛盾。所以定义域为 $I$。
\end{solution}


\Needspace{8\baselineskip}
\section{拓扑空间}\label{sec:yosida-00-02}
\noindent\textit{原书位置：中译本印刷第 3--14 页；英文版第 3--15 页。}
% !TeX encoding = UTF-8
本节以开集组织邻近、连续与紧性，随后证明函数逼近和类别论中常用的基础定理。除特别说明外，拓扑空间均假设为 Hausdorff 空间；个别结论会明确指出不需要这一假设。紧性始终指开覆盖意义的紧性。

\subsection{开集、闭集、邻域与连续映射}
\begin{definition}[拓扑与子空间拓扑]\label{def:y0-topology}
集合 $X$ 上的拓扑是子集族 $\tau\subseteq\mathcal P(X)$，满足 $\varnothing,X\in\tau$，任意族成员的并属于 $\tau$，有限个成员的交属于 $\tau$。$\tau$ 的成员称为开集，补集为开集的集合称为闭集。子集 $A\subseteq X$ 的子空间拓扑为 $\{A\cap U:U\in\tau\}$。

若 $N$ 含有一个包含 $x$ 的开集，则称 $N$ 为 $x$ 的邻域。若任意两个不同点都有互不相交的开邻域，则空间称为 Hausdorff 空间。
\end{definition}

\begin{definition}[闭包、内部、边界与聚点]\label{def:y0-closure}
集合 $A$ 的闭包 $\overline A$ 是包含 $A$ 的全部闭集的交；内部 $A^\circ$ 是包含于 $A$ 的全部开集的并；边界为 $\partial A=\overline A\setminus A^\circ$。若 $\overline A=X$，称 $A$ 在 $X$ 中稠密。若每个 $x$ 的邻域都与 $A\setminus\{x\}$ 相交，称 $x$ 为 $A$ 的聚点。
\end{definition}

\begin{proposition}[闭包的邻域刻画]\label{prop:y0-closure}
$x\in\overline A$ 当且仅当 $x$ 的每个开邻域都与 $A$ 相交。因此 $A$ 闭当且仅当它包含自己的全部聚点，且
\[
 \overline A=X\setminus(X\setminus A)^\circ.
\]
\end{proposition}
\begin{proof}
若某开邻域 $U\ni x$ 与 $A$ 不交，则闭集 $X\setminus U$ 包含 $A$，而不含 $x$，所以 $x\notin\overline A$。反之，若 $x\notin\overline A$，则 $X\setminus\overline A$ 是所求开邻域。对 $x\notin A$，与 $A$ 相交等价于与 $A\setminus\{x\}$ 相交，故闭集的刻画随之成立。最后一个等式是邻域刻画的补集形式。
\end{proof}

\begin{definition}[拓扑基、邻域基与连续]\label{def:y0-continuity}
开集族 $\mathcal B$ 若使每个开集都是其成员的并，称为拓扑基。点 $x$ 的邻域族 $\mathcal B_x$ 若使每个 $x$ 的邻域都包含一个 $\mathcal B_x$ 的成员，称为 $x$ 的邻域基。每点有可数邻域基称为第一可数；有可数拓扑基称为第二可数。

映射 $f:X\to Y$ 在 $x$ 连续，指对每个 $f(x)$ 的邻域 $V$，存在 $x$ 的邻域 $U$ 使 $f(U)\subseteq V$；在每点连续称为连续。双射若其本身和逆映射都连续，称为同胚。
\end{definition}

\begin{proposition}[连续性的等价条件]\label{prop:y0-continuity}
映射 $f:X\to Y$ 连续当且仅当每个开集的原像为开集，也当且仅当每个闭集的原像为闭集。
\end{proposition}
\begin{proof}
若 $f$ 连续，对开集 $V\subseteq Y$ 及 $x\in f^{-1}(V)$，连续性给出含 $x$ 的开集 $U\subseteq f^{-1}(V)$，所以原像开。反之，给定邻域 $V\ni f(x)$，在其中取开邻域 $W$，则 $f^{-1}(W)$ 为 $x$ 的开邻域，且像包含于 $V$。闭集条件与开集条件由取补集互相转换。
\end{proof}

\begin{definition}[度量与收敛]\label{def:y0-metric}
函数 $d:X\times X\to[0,\infty)$ 若满足 $d(x,y)=0$ 当且仅当 $x=y$、$d(x,y)=d(y,x)$ 和三角不等式 $d(x,z)\le d(x,y)+d(y,z)$，称为度量。开球 $B(x,r)=\{y:d(x,y)<r\}$ 构成拓扑基，所定义的拓扑称为度量拓扑。

在拓扑空间中，$x_n\to x$ 指每个 $x$ 的邻域都包含数列的全部充分靠后项。在度量空间中，这等价于 $d(x_n,x)\to0$。
\end{definition}

度量拓扑是 Hausdorff 的：不同点之间距离为 $r>0$，半径 $r/3$ 的两球不交。一般 Hausdorff 空间中数列极限也唯一，否则两个不交的极限邻域都必须含有充分靠后的项。度量空间是第一可数的，因为 $B(x,1/n)$ 给出可数邻域基。

\begin{proposition}[第一可数空间中的数列判据]\label{prop:y0-sequential}
若 $X$ 第一可数，则 $x\in\overline A$ 当且仅当存在 $A$ 中的数列趋于 $x$；映射 $f:X\to Y$ 连续当且仅当它保持所有数列的收敛。
\end{proposition}
\begin{proof}
将 $x$ 的可数邻域基替换为有限交，使它递减，记为 $U_n$。若 $x\in\overline A$，取 $x_n\in A\cap U_n$，便有 $x_n\to x$。反向由闭包的邻域刻画成立。连续映射显然保持收敛；若 $f$ 在 $x$ 不连续，则有开邻域 $V\ni f(x)$ 使每个 $U_n$ 都含有点 $x_n$ 满足 $f(x_n)\notin V$，这给出反例数列。
\end{proof}

一般拓扑空间中数列未必刻画闭包或连续性。例如在不可数集的可数补拓扑中，每个收敛数列最终恒等于其极限，但所有不可数子集都稠密。此例不满足 Hausdorff 条件，仅用来说明不能省略数列判据的假设。

\subsection{紧性、局部紧性与单点紧化}
\begin{definition}[紧性与相对紧性]\label{def:y0-compact}
若子集 $K\subseteq X$ 的每个开覆盖都有有限子覆盖，称 $K$ 紧。若 $\overline A$ 紧，称 $A$ 相对紧。若每个点都具有紧邻域，称空间局部紧。

集合族具有有限交性质，指每个有限子族的交都非空，其中空族的交约定为整个底集。
\end{definition}

\begin{proposition}[紧性的基本性质]\label{prop:y0-compact}
连续像保持紧性；紧空间的闭子集紧；Hausdorff 空间的紧子集闭；紧性等价于每个具有有限交性质的闭集族的交非空。紧空间到 Hausdorff 空间的连续双射为同胚。
\end{proposition}
\begin{proof}
像空间的开覆盖取原像，再用有限子覆盖，即得第一条。覆盖闭子集时加上其补集，得到紧空间的覆盖，即得第二条。若 $x\notin K$ 且 $K$ 紧，对每个 $y\in K$ 取不交开邻域 $U_y\ni y$、$V_y\ni x$。有限个 $U_{y_j}$ 覆盖 $K$，则 $\bigcap_jV_{y_j}$ 与 $K$ 不交，故 $X\setminus K$ 开。有限交刻画由开覆盖取补集得到。最后，连续双射将闭集送到紧集，再送到闭集，故其逆映射连续。
\end{proof}

\begin{proposition}[紧 Hausdorff 空间的正规性]\label{prop:y0-normal}
紧 Hausdorff 空间是正规空间，即任意两个不交闭集都有互不相交的开邻域。它也是正则空间，即点与不含它的闭集可被不交开集分离。
\end{proposition}
\begin{proof}
设不交闭集 $A,B$，均紧。固定 $a\in A$，对每个 $b\in B$ 取不交开邻域 $U_{a,b}\ni a$、$V_{a,b}\ni b$。以有限个 $V_{a,b_j}$ 覆盖 $B$，令 $U_a=\bigcap_jU_{a,b_j}$、$V_a=\bigcup_jV_{a,b_j}$。则 $U_a\cap V_a=\varnothing$ 且 $B\subseteq V_a$。再以有限个 $U_{a_i}$ 覆盖 $A$，取 $U=\bigcup_iU_{a_i}$、$V=\bigcap_iV_{a_i}$，便有 $A\subseteq U$、$B\subseteq V$、$U\cap V=\varnothing$。单点在 Hausdorff 空间中闭，故正规性也给出正则性。
\end{proof}

\begin{theorem}[单点紧化]\label{thm:y0-one-point}
设 $X$ 局部紧且 Hausdorff，添加新点 $\infty$，令 $X^*=X\cup\{\infty\}$。规定 $X$ 原有开集仍开，含 $\infty$ 的集合 $W$ 开当且仅当 $X^*\setminus W$ 是 $X$ 的紧子集。则 $X^*$ 为紧 Hausdorff 空间，且诱导于 $X$ 的拓扑不变。若 $X$ 非紧，则 $X$ 在 $X^*$ 中稠密；若 $X$ 紧，则 $\infty$ 为孤立点。
\end{theorem}
\begin{proof}
紧子集闭；含 $\infty$ 的开集取有限交时，补集是有限个紧集的并，仍紧。取任意并时，若其中某成员含 $\infty$，并集的补集是该成员补集的闭子集，故紧。与原有开集取交也保持开性。这证明拓扑公理及子空间拓扑不变。

任意覆盖 $X^*$ 的开族中，取一个含 $\infty$ 的成员，其补集紧，可由有限个其余成员覆盖，故 $X^*$ 紧。$X$ 中两点沿用原来的分离。对 $x\in X$，取紧邻域 $K$ 及开集 $U$ 使 $x\in U\subseteq K$，则 $U$ 和 $X^*\setminus K$ 分离 $x$ 与 $\infty$。若 $X$ 非紧，$\infty$ 的任何邻域都与 $X$ 相交；若 $X$ 紧，$\{\infty\}$ 开。
\end{proof}

\subsection{乘积拓扑与 Tychonoff 定理}
\begin{definition}[乘积拓扑]\label{def:y0-product-topology}
对空间族 $(X_i)_{i\in I}$，乘积 $X=\prod_iX_i$ 的乘积拓扑以集合 $\bigcap_{i\in F}\pi_i^{-1}(U_i)$ 为基，其中 $F\subseteq I$ 有限，各 $U_i\subseteq X_i$ 开，$\pi_i$ 为坐标投影。它是使所有投影连续的最弱拓扑。若允许无限多个坐标同时受限，得到的箱拓扑通常更强。
\end{definition}

以下证明使用第 \ref{thm:y0-zorn} 定理，因而在选择公理下进行。为使任意多个坐标的极限相容，先引入滤子。

\begin{lemma}[超滤子与紧性]\label{lem:y0-ultrafilter}
非空集合 $X$ 上的滤子是子集族 $\mathcal F$，满足 $X\in\mathcal F$、$\varnothing\notin\mathcal F$，对有限交封闭，并且包含成员的所有超集。极大的滤子称为超滤子。每个滤子可扩充为超滤子；超滤子 $\mathcal U$ 对任意 $A\subseteq X$ 恰包含 $A$、$X\setminus A$ 之一。

称滤子收敛于 $x$，若它包含 $x$ 的所有邻域。空间紧当且仅当其上每个超滤子都收敛于至少一点；Hausdorff 空间中的滤子极限至多一个。
\end{lemma}
\begin{proof}
所有包含固定滤子的滤子按包含排序，链的并仍为滤子。Zorn 引理给出极大者。若 $A\notin\mathcal U$，将 $A$ 加入不能生成真滤子，故存在 $F\in\mathcal U$ 使 $F\cap A=\varnothing$；于是 $X\setminus A\in\mathcal U$。两者不能同时属于滤子。

若空间紧，闭集族 $\{\overline F:F\in\mathcal U\}$ 具有有限交性质，故取 $x$ 在其交中。若邻域 $V\ni x$ 不在 $\mathcal U$ 中，取开邻域 $W\subseteq V$，则 $X\setminus W\in\mathcal U$，与 $x\in\overline{X\setminus W}=X\setminus W$ 矛盾。因此超滤子收敛于 $x$。反之，具有有限交性质的闭集族生成滤子，将其扩充为超滤子并取极限 $x$。若某个原闭集 $F$ 不含 $x$，其补集为 $x$ 的邻域，与 $F$ 同时在滤子中，矛盾，所以 $x$ 属于全部原闭集。Hausdorff 情形的唯一性由不交邻域不能同时属于滤子得到。
\end{proof}

\begin{theorem}[Tychonoff 定理]\label{thm:y0-tychonoff}
任意族紧 Hausdorff 空间的乘积，在乘积拓扑下仍为紧 Hausdorff 空间。
\end{theorem}
\begin{proof}
若乘积为空，它紧。否则取乘积上的超滤子 $\mathcal U$。对每个 $i$，族 $\mathcal U_i=\{A\subseteq X_i:\pi_i^{-1}(A)\in\mathcal U\}$ 为 $X_i$ 上的超滤子，由紧性它收敛于 $x_i$。令 $x=(x_i)_{i\in I}$。$x$ 的每个基本邻域是有限个 $\pi_i^{-1}(U_i)$ 的交，各项都在 $\mathcal U$ 中，故交也在其中。这说明 $\mathcal U\to x$，由引理得乘积紧。不同乘积点在某一坐标不同，该坐标中的不交开邻域的原像分离这两点，故乘积 Hausdorff。
\end{proof}

\subsection{Urysohn 引理与连续截断函数}
\begin{theorem}[Urysohn 引理]\label{thm:y0-urysohn}
设 $X$ 为正规 Hausdorff 空间，$A,B\subseteq X$ 为不交闭集。则存在连续函数 $f:X\to[0,1]$，使 $f|_A=0$、$f|_B=1$。
\end{theorem}
\begin{proof}
正规性给出缩小邻域的性质：若闭集 $F\subseteq O$ 且 $O$ 开，则存在开集 $V$ 使 $F\subseteq V\subseteq\overline V\subseteq O$。确实，将 $F$ 与 $X\setminus O$ 用不交开集分离即可。

令 $D$ 为 $[0,1]$ 中的二进有理数。取开集 $U_1=X\setminus B$，再取 $U_0$ 使 $A\subseteq U_0\subseteq\overline U_0\subseteq U_1$。逐层在相邻已定义指标 $r<s$ 的中点插入开集，使
\[
 \overline U_r\subseteq U_{(r+s)/2}
 \subseteq\overline U_{(r+s)/2}\subseteq U_s.
\]
从而对任意 $r<s$ 都有 $\overline U_r\subseteq U_s$。定义
\[
 f(x)=\inf\bigl(\{r\in D:x\in U_r\}\cup\{1\}\bigr).
\]
显然 $f=0$ 于 $A$，$f=1$ 于 $B$。对 $0<a\le1$，有 $\{f<a\}=\bigcup_{r<a}U_r$。对 $0\le a<1$，有
\[
 \{f>a\}=\bigcup_{\substack{r\in D\\r>a}}
                     (X\setminus\overline U_r).
\]
后一等式中，若 $x\notin\overline U_r$，则 $x$ 不属于任何 $U_s$（$s\le r$），所以 $f(x)\ge r>a$；若 $f(x)>a$，在 $a$ 与 $f(x)$ 之间取 $r<s$，则 $x\notin U_s$，从而 $x\notin\overline U_r$。这些原像均开，证明连续性。
\end{proof}

\begin{corollary}[局部紧空间中的截断函数]\label{cor:y0-cutoff}
若 $X$ 局部紧且 Hausdorff，$K\subseteq O\subseteq X$，其中 $K$ 紧、$O$ 开，则有 $\varphi\in C_c(X)$ 满足 $0\le\varphi\le1$、$\varphi|_K=1$、$\operatorname{supp}\varphi\subseteq O$。这里 $C_c(X)$ 表示连续且具有紧支集的函数。
\end{corollary}
\begin{proof}
在紧 Hausdorff 空间 $X^*$ 中，$K$ 闭而 $O$ 开，由正规性取开集 $V$ 使 $K\subseteq V\subseteq\overline V\subseteq O$。对不交闭集 $K$、$X^*\setminus V$ 应用 Urysohn 引理，得到在 $K$ 上为 $1$、在后者上为 $0$ 的函数。限制到 $X$ 后，其支集包含于紧集 $\overline V\subseteq O$，即得结论。
\end{proof}

\subsection{Weierstrass 与 Stone--Weierstrass 逼近}
\begin{theorem}[Bernstein 多项式逼近]\label{thm:y0-bernstein}
对 $f\in C([0,1];\mathbb K)$，其中 $\mathbb K=\R$ 或 $\mathbb C$，多项式
\[
 B_nf(x)=\sum_{k=0}^nf(k/n)\binom nk x^k(1-x)^{n-k}
\]
一致收敛于 $f$。因此任意闭有限区间上的连续函数均可由多项式一致逼近。
\end{theorem}
\begin{proof}
记 $p_{n,k}(x)=\binom nkx^k(1-x)^{n-k}$。二项式公式及其一、二阶导数给出
\[
 \sum_kp_{n,k}=1,\qquad
 \sum_k(k/n-x)^2p_{n,k}=\frac{x(1-x)}n\le\frac1{4n}.
\]
设 $M=\|f\|_\infty$。由紧区间上的一致连续性，给定 $\varepsilon>0$ 可取 $\delta>0$ 使 $|s-t|<\delta$ 时 $|f(s)-f(t)|<\varepsilon$。将求和按 $|k/n-x|<\delta$ 与其补集分开，得
\[
 |B_nf(x)-f(x)|
 \le\varepsilon+2M\delta^{-2}\sum_k(k/n-x)^2p_{n,k}
 \le\varepsilon+\frac{M}{2n\delta^2}.
\]
估计对 $x$ 一致，故结论成立。一般区间由仿射换元得到。这里的一致连续性可由紧性证明：在每点取控制函数振荡的小邻域，有限子覆盖再给出统一的距离界限。
\end{proof}

\begin{theorem}[实 Stone--Weierstrass 定理]\label{thm:y0-stone-real}
设 $X$ 紧 Hausdorff，$\mathcal A\subseteq C(X;\R)$ 为含常数函数且分离点的子代数；分离点指对 $x\ne y$ 存在 $a\in\mathcal A$ 使 $a(x)\ne a(y)$。则 $\mathcal A$ 在一致范数 $\|f\|_\infty=\sup_X|f|$ 下稠密于 $C(X;\R)$。
\end{theorem}
\begin{proof}
设 $\mathcal B$ 为 $\mathcal A$ 的一致闭包，它仍为含常数的代数。对 $g\in\mathcal B$，在包含其值域的有限区间上以多项式逼近绝对值函数，得 $|g|\in\mathcal B$。因此
\[
 \max(g,h)=\tfrac12(g+h+|g-h|),\qquad
 \min(g,h)=\tfrac12(g+h-|g-h|)
\]
也在 $\mathcal B$ 中。

给定 $f\in C(X;\R)$ 与 $\varepsilon>0$。对任意 $x,y\in X$，可找 $a_{x,y}\in\mathcal A$ 满足 $a_{x,y}(x)=f(x)$、$a_{x,y}(y)=f(y)$：若 $x\ne y$，对一个分离它们的函数作仿射变换；若 $x=y$，取常数函数。固定 $x$，开集 $\{t:a_{x,y}(t)>f(t)-\varepsilon\}$ 随 $y$ 覆盖 $X$。取有限子覆盖并令 $g_x$ 为对应函数的最大值，则 $g_x\in\mathcal B$、$g_x>f-\varepsilon$ 于 $X$，且 $g_x(x)=f(x)$。

开集 $\{t:g_x(t)<f(t)+\varepsilon\}$ 随 $x$ 覆盖 $X$。再取有限子覆盖，令 $g$ 为对应 $g_x$ 的最小值，便有 $f-\varepsilon<g<f+\varepsilon$ 于 $X$，且 $g\in\mathcal B$。由闭性 $f\in\mathcal B$，完成证明。
\end{proof}

\begin{corollary}[格形式与复数形式]\label{cor:y0-stone-complex}
实数情形中，可将“子代数”替换为“对有限最大、最小运算封闭的线性子空间”，仍要求含常数及分离点。复数情形中，含常数、分离点且对复共轭封闭的子代数 $\mathcal A\subseteq C(X;\mathbb C)$ 在一致范数下稠密。
\end{corollary}
\begin{proof}
格形式直接使用上一证明的两次有限覆盖构造。复数情形中，实部和虚部都属于 $\mathcal A$，故其中的实值函数构成含常数的实代数，并分离点：一个复值函数在两点不同，其实部或虚部必不同。对任意复值连续函数分别逼近实部和虚部，再组合即可。
\end{proof}

\begin{example}[三角多项式的一致稠密性]
在单位圆上，函数 $z\mapsto z^k$（$k\in\mathbb Z$）的有限复线性组合构成含常数、分离点且对共轭封闭的代数。因此每个连续 $2\pi$ 周期函数都可由 $\sum_{k=-N}^Nc_ke^{ikt}$ 一致逼近。这里只断言存在逼近的三角多项式，不断言原函数的 Fourier 部分和必一致收敛。
\end{example}

\subsection{完备性与 Baire 类别定理}
\begin{definition}[Cauchy 数列、完备性与类别]\label{def:y0-baire}
度量空间中，数列 $(x_n)$ 为 Cauchy 数列，指对任意 $\varepsilon>0$，存在 $N$ 使 $m,n\ge N$ 时 $d(x_m,x_n)<\varepsilon$。每个 Cauchy 数列都在空间内收敛，称空间完备。

若 $(\overline A)^\circ=\varnothing$，称 $A$ 无处稠密；若 $A$ 是可数个无处稠密集的并，称 $A$ 为第一纲集或贫集；否则称为第二纲集。若可数个稠密开集的交仍稠密，称空间为 Baire 空间。
\end{definition}

\begin{theorem}[Baire 类别定理]\label{thm:y0-baire}
每个完备度量空间和每个局部紧 Hausdorff 空间都是 Baire 空间。特别地，非空完备度量空间或非空紧 Hausdorff 空间不能是自身中的第一纲集。
\end{theorem}
\begin{proof}
设 $G_n$ 为稠密开集，$O$ 为非空开集。在完备度量空间中，递归取非空闭球 $K_n=\overline B(x_n,r_n)$，使 $K_1\subseteq O\cap G_1$、$K_{n+1}\subseteq B(x_n,r_n)\cap G_{n+1}$，并使 $0<r_n<2^{-n}$。稠密性保证每一步可进行。中心为 Cauchy 数列，由完备性收敛于 $x$；闭球嵌套给出 $x\in K_n$ 对所有 $n$ 成立，故 $x\in O\cap\bigcap_nG_n$。

在局部紧 Hausdorff 空间中，每个非空开集包含一个非空开集 $V$，使 $\overline V$ 紧且包含于原开集；这由单点紧化空间的正则性得到。于是递归取 $V_n$ 使 $\overline V_1\subseteq O\cap G_1$、$\overline V_{n+1}\subseteq V_n\cap G_{n+1}$，且各闭包紧。闭集 $\overline V_n$ 在 $\overline V_1$ 中具有有限交性质，交非空，交中的点即属于 $O\cap\bigcap_nG_n$。最后，对无处稠密集的闭包取补集，便得到第一纲表述。
\end{proof}

\begin{proposition}[连续函数逐点极限的不连续点]\label{prop:y0-baire-one}
设 $f_n:X\to\R$ 连续，且在每点收敛于有限值 $f$。则 $f$ 的不连续点集为第一纲集。若 $X$ 为 Baire 空间，$f$ 的连续点集稠密。
\end{proposition}
\begin{proof}
固定 $\varepsilon>0$，令
\[
 F_N=\bigcap_{m,n\ge N}\{x:|f_m(x)-f_n(x)|\le\varepsilon/4\}.
\]
各 $F_N$ 闭，且由逐点 Cauchy 性 $X=\bigcup_NF_N$。在 $F_N$ 上取极限得 $|f-f_N|\le\varepsilon/4$。若 $x\in F_N^\circ$，由 $f_N$ 的连续性，可取 $x$ 的开邻域 $U\subseteq F_N^\circ$，使 $f_N$ 在 $U$ 上的振荡小于 $\varepsilon/2$；于是 $f$ 在 $U$ 上的振荡小于 $\varepsilon$。

令 $G_\varepsilon$ 为所有使 $f$ 的振荡小于 $\varepsilon$ 的开集的并。则 $G_\varepsilon$ 开，且
\[
 X\setminus G_\varepsilon\subseteq
          \bigcup_N(F_N\setminus F_N^\circ).
\]
右侧为第一纲集，因为闭集 $F_N$ 的边界无处稠密。$f$ 在 $x$ 连续当且仅当 $x\in\bigcap_{k\ge1}G_{1/k}$，故不连续点集第一纲。在 Baire 空间中，第一纲集不能包含非空开集，因而其补集稠密。
\end{proof}

\subsection{度量空间中的紧性判据}
\begin{definition}[全有界]\label{def:y0-total-bounded}
度量空间 $M$ 若对每个 $\varepsilon>0$ 都能用有限个半径 $\varepsilon$ 的球覆盖，称为全有界。对子集使用限制的度量。全有界蕴含有界，但无限维空间中有界未必全有界。
\end{definition}

\begin{theorem}[度量紧性的等价条件]\label{thm:y0-metric-compact}
度量空间 $M$ 的下列条件等价：紧；每个数列有收敛子列且极限在 $M$ 中；完备且全有界。因此完备度量空间 $X$ 中的子集 $A$ 相对紧，当且仅当 $A$ 全有界。
\end{theorem}
\begin{proof}
紧空间全有界，因为固定半径的全部开球构成覆盖。紧度量空间也完备：对 Cauchy 数列，闭集 $\overline{\{x_n:n\ge N\}}$ 嵌套且非空，紧性给出公共点 $x$，Cauchy 性随即给出 $x_n\to x$。

完备且全有界时，依次用半径 $2^{-k}$ 的有限球覆盖，从当前无限多个数列项中保留落在同一球内的无限子列。取对角子列，得到 Cauchy 子列，再由完备性得到收敛。反之，数列紧性蕴含全有界：若不全有界，递归选取两两距离至少某个 $\varepsilon>0$ 的点，便没有收敛子列。

最后由数列紧性推出开覆盖紧性。给定开覆盖，必有 $\delta>0$ 使每个球 $B(x,\delta)$ 都包含于某一覆盖成员。否则可取 $x_n$ 使 $B(x_n,1/n)$ 不包含于任何成员；取收敛子列 $x_{n_j}\to x$，再取覆盖成员 $U\ni x$ 和 $r>0$ 使 $B(x,r)\subseteq U$，则充分大时 $B(x_{n_j},1/n_j)\subseteq U$，矛盾。用全有界性取有限个半径 $\delta/2$ 的球，再为其中心各选一个包含半径 $\delta$ 球的覆盖成员，即得有限子覆盖。

若 $X$ 完备，$\overline A$ 为完备子空间；$A$ 全有界时，放大覆盖球半径即知其闭包也全有界，故闭包紧。反向由紧集全有界及其子集仍全有界成立。
\end{proof}

\begin{example}[有界不足以推出紧]
在 $\ell^2$ 中，标准单位向量 $e_n$ 均在闭单位球内，但 $\|e_n-e_m\|_2=\sqrt2$（$n\ne m$），所以闭单位球不全有界，也不紧。这与 $\R^n$ 中“闭且有界等价于紧”形成区别；后者由有限网格覆盖给出全有界性，再结合完备性证明。
\end{example}

\begin{exercise}[紧空间上的连续函数空间]
设 $K$ 紧 Hausdorff。证明 $C(K;\mathbb K)$ 在一致范数下完备。
\end{exercise}
\begin{solution}
一致 Cauchy 数列 $(f_n)$ 在每点为标量 Cauchy 数列，令 $f(x)=\lim_nf_n(x)$。将 $\|f_n-f_m\|_\infty<\varepsilon$ 中的 $m$ 取极限，得 $|f_n(x)-f(x)|\le\varepsilon$ 对所有 $x$ 成立，故一致收敛。固定 $x$，先使 $\|f-f_n\|_\infty$ 小，再以 $f_n$ 的连续性控制邻域中的变化，由三角不等式知 $f$ 连续。连续函数在紧空间上的像紧，因而有界，故 $f\in C(K;\mathbb K)$。
\end{solution}

\begin{exercise}[复 Stone--Weierstrass 中的共轭条件]
在单位圆 $\mathbb T$ 上，令 $\mathcal A$ 为仅含非负幂 $z^k$ 的多项式的限制。证明它含常数且分离点，却不能一致逼近 $\overline z$。
\end{exercise}
\begin{solution}
函数 $z$ 本身分离点。对任意 $p\in\mathcal A$，有 $\int_0^{2\pi}p(e^{it})e^{it}\,\dd t=0$，因为每个指数的频率均为正整数。若 $p_n\to\overline z$ 一致，则左侧趋于 $\int_0^{2\pi}e^{-it}e^{it}\,\dd t=2\pi$，矛盾。因此共轭封闭条件不能直接删去。
\end{solution}


\Needspace{8\baselineskip}
\section{测度空间}\label{sec:yosida-00-03}
\noindent\textit{原书位置：中译本印刷第 14--20 页；英文版第 15--20 页。}
% !TeX encoding = UTF-8
本节采用现代的 $\sigma$-代数记号，先从简单函数构造积分，再证明极限交换与乘积积分定理。原书将 $\sigma$-有限性列入主要工作设定；本稿将它单独标出，使读者能看清各结论实际需要什么条件。除涉及乘积测度等处另加说明外，一般测度空间不预设拓扑。

\subsection{可测集、测度与几乎处处}
\begin{definition}[$\sigma$-代数与测度空间]\label{def:y0-measure}
集合 $S$ 上的 $\sigma$-代数 $\Sigma$ 是包含 $S$、对补集和可数并封闭的子集族；其成员称为可测集。映射 $\mu:\Sigma\to[0,\infty]$ 若满足 $\mu(\varnothing)=0$，且对两两不交的 $E_n\in\Sigma$ 有
\[
 \mu\left(\bigcup_{n\ge1}E_n\right)=\sum_{n\ge1}\mu(E_n),
\]
称为测度，$(S,\Sigma,\mu)$ 称为测度空间。若 $\mu(S)<\infty$，称测度有限；若 $S=\bigcup_nS_n$ 且每个 $\mu(S_n)<\infty$，称测度 $\sigma$-有限。

若性质在某个可测零测集之外成立，称它几乎处处成立，简记为 a.e.。若每个可测零测集的所有子集均可测，称测度空间完备。
\end{definition}

原书称含 $S$ 且对补集、可数并封闭的族为“$\sigma$-ring”；按现代通常术语，这是 $\sigma$-代数。现代 $\sigma$-环只要求对差和可数并封闭，不一定包含 $S$。这个区别在非 $\sigma$-紧空间上的紧集生成测度中尤其需要注意。

\begin{proposition}[测度的基本极限性质]\label{prop:y0-measure-continuity}
测度单调且可数次可加。若 $E_n\uparrow E$，则 $\mu(E_n)\uparrow\mu(E)$；若 $E_n\downarrow E$ 且 $\mu(E_1)<\infty$，则 $\mu(E_n)\downarrow\mu(E)$。
\end{proposition}
\begin{proof}
将 $B$ 写为 $A$ 与 $B\setminus A$ 的不交并得单调性。一般可数并可按首次出现的指标分成不交差集，得次可加性。对递增族，将并写为 $E_1$ 与 $E_n\setminus E_{n-1}$ 的不交并，级数部分和即 $\mu(E_n)$。递减情形对 $E_1\setminus E_n$ 使用递增结论，再从有限数 $\mu(E_1)$ 中相减。有限性不能省略，例如计数测度下 $\{n,n+1,\ldots\}\downarrow\varnothing$，每项测度却为无穷。
\end{proof}

\begin{proposition}[测度的完备化]\label{prop:y0-completion-measure}
令 $\widehat\Sigma$ 为所有形如 $E\cup N$ 的集合，其中 $E\in\Sigma$，$N$ 包含于某个 $\Sigma$ 中的零测集。则 $\widehat\Sigma$ 为 $\sigma$-代数，且 $\widehat\mu(E\cup N)=\mu(E)$ 给出唯一延拓 $\mu$ 的完备测度。
\end{proposition}
\begin{proof}
每个新集合 $A$ 都夹在两个原可测集 $E\subseteq A\subseteq F$ 之间，且 $\mu(F\setminus E)=0$。取补集或可数并仍可找到这样的上下夹集，证明 $\widehat\Sigma$ 为 $\sigma$-代数。不同夹集的测度相同，因为它们只相差零测集，所以定义良定。对不交的新集合 $A_n$，选择可测下夹集 $E_n\subseteq A_n$，它们也不交，且 $\bigcup_nA_n$ 与 $\bigcup_nE_n$ 只相差零测集。可数可加性因此继承。完备性及唯一性由零测集的任何子集必须有测度零得到。
\end{proof}

\subsection{可测函数与 Egorov 定理}
\begin{definition}[可测函数与简单函数]\label{def:y0-measurable-function}
函数 $f:S\to[-\infty,\infty]$ 称为可测，若 $\{s:f(s)>a\}\in\Sigma$ 对每个 $a\in\R$ 成立。有限复值函数可测，指其实部与虚部可测，等价于每个复平面开集的原像可测。取有限个值的可测函数称为简单函数，可写作 $\sum_{j=1}^ma_j\mathbf1_{E_j}$，其中 $E_j$ 两两不交且可测。
\end{definition}

可测函数的有限和、乘积、绝对值仍可测；可数上确界、下确界、上极限及下极限也可测。例如 $\{\sup_nf_n>a\}=\bigcup_n\{f_n>a\}$，下确界与极限由类似可数运算表达。每个非负可测函数都有递增的非负简单函数逼近；可取
\[
 s_n(s)=2^{-n}\left\lfloor2^n\min(f(s),n)\right\rfloor,
 \qquad s_n\uparrow f,
\]
其中在 $f=\infty$ 处规定截断值为 $n$。每个取值区间的原像可测，故 $s_n$ 可测。

\begin{theorem}[Egorov 定理]\label{thm:y0-egorov}
设 $E\in\Sigma$、$\mu(E)<\infty$，$f_n$、$f$ 为可测函数，在 $E$ 上几乎处处有限，且 $f_n\to f$ 几乎处处。则对每个 $\eta>0$，存在可测集 $F\subseteq E$，使 $\mu(E\setminus F)<\eta$，且 $f_n\to f$ 在 $F$ 上一致。
\end{theorem}
\begin{proof}
先从 $E$ 中去掉一个可测零测集，使全部函数有限并逐点收敛。对正整数 $k,N$，令
\[
 A_{k,N}=\bigcup_{n\ge N}\{s\in E:|f_n(s)-f(s)|>1/k\}.
\]
固定 $k$ 时，$A_{k,N}\downarrow\varnothing$；由有限测度下的上连续性，可取 $N_k$ 使 $\mu(A_{k,N_k})<\eta2^{-k}$。令 $F=E\setminus\bigcup_kA_{k,N_k}$，则损失测度小于 $\eta$，且 $n\ge N_k$ 时 $\sup_F|f_n-f|\le1/k$。这正是一致收敛。
\end{proof}

\begin{example}[有限测度条件的作用]
在 $\R$ 的 Lebesgue 测度下，$f_n=\mathbf1_{[n,\infty)}$ 逐点趋于零；但从 $\R$ 去掉任意有限测度集后，剩余集合仍与每个 $[n,\infty)$ 相交，故误差上确界始终为 $1$。不能将 Egorov 定理中的有限测度条件直接换成 $\sigma$-有限。
\end{example}

\subsection{积分的构造与单调收敛}
\begin{definition}[Lebesgue 积分]\label{def:y0-integral}
对非负简单函数 $s=\sum_ja_j\mathbf1_{E_j}$，定义 $\int s\,\dd\mu=\sum_ja_j\mu(E_j)$，约定 $0\cdot\infty=0$。对非负可测函数 $f$，定义
\[
 \int f\,\dd\mu=\sup\left\{\int s\,\dd\mu:
                      0\le s\le f,\ s\text{ 为简单函数}\right\}.
\]
实值函数写为 $f=f^+-f^-$，其中 $f^+=\max(f,0)$、$f^-=\max(-f,0)$。若 $\int|f|\,\dd\mu<\infty$，称 $f$ 可积，并定义 $\int f=\int f^+-\int f^-$。复值函数可积也指 $\int|f|<\infty$，积分分别取实部和虚部的积分。局部可积、$L^p$ 等概念将在后续章节使用。
\end{definition}

对简单函数，不同表示可细分到同一不交分割，因此积分良定、单调、可加且正齐次。对于一般函数，先证明下一定理，再从简单函数逼近继承这些性质。

\begin{theorem}[单调收敛定理]\label{thm:y0-monotone}
若 $0\le f_n\uparrow f$ 为可测函数，则 $\int f_n\,\dd\mu\uparrow\int f\,\dd\mu$，允许积分无穷。
\end{theorem}
\begin{proof}
由单调性，左侧极限 $L\le\int f$。取简单函数 $0\le s\le f$ 和 $0<c<1$。令 $E_n=\{f_n\ge cs\}$。这些集合递增，其并包含 $\{s>0\}$：若 $s(x)>0$，则 $f_n(x)\uparrow f(x)\ge s(x)>cs(x)$。所以由测度的下连续性和简单函数公式，$\int s\mathbf1_{E_n}\to\int s$。于是
\[
 \int f_n\ge c\int s\mathbf1_{E_n},\qquad
 L\ge c\int s.
\]
令 $c\uparrow1$，再对 $s$ 取上确界，得 $L\ge\int f$。
\end{proof}

非负函数的积分可加性现在可由各自的递增简单逼近和单调收敛得出；有限实值及复值可积函数的线性性由正负部或实虚部的分解得出。特别地，可积性在加法下保持，且
\[
 \left|\int f\,\dd\mu\right|\le\int|f|\,\dd\mu.
\]
复数情形可选模为 $1$ 的标量 $\alpha$，使 $\alpha\int f$ 为非负实数，再用 $\operatorname{Re}(\alpha f)\le|f|$ 得到估计。若 $f\ge0$，则 $\int f=0$ 当且仅当 $f=0$ a.e.：因为 $\mu(\{f\ge1/n\})\le n\int f$，而 $\{f>0\}$ 是这些集合的并。

\begin{proposition}[与简单函数逼近定义的相容性]\label{prop:y0-simple-l1}
每个可积函数 $f$ 都有可积简单函数 $s_n$ 满足 $\int|s_n-f|\to0$。反之，若可测有限函数 $f$ 是可积简单函数 $s_n$ 的几乎处处极限，且 $(s_n)$ 在积分距离 $\int|s_n-s_m|$ 下 Cauchy，则 $f$ 可积，$\int|s_n-f|\to0$，并且 $\int s_n\to\int f$。
\end{proposition}
\begin{proof}
第一条先对 $f\ge0$ 使用递增简单逼近，由单调收敛及 $\int f<\infty$ 得 $\int(f-s_n)\to0$；再分别处理实虚部和正负部。反向将在下面 Fatou 引理得到后使用它：固定 $n$，
\[
 \int|s_n-f|\le\liminf_{m\to\infty}\int|s_n-s_m|.
\]
右侧在 $n$ 充分大时任意小，故 $f$ 与某个可积 $s_n$ 的差可积，从而 $f$ 可积，并得积分距离趋零。积分估计又给出积分值的收敛。这说明本稿的非负积分构造与原书按简单函数积分距离定义的方式相容。
\end{proof}

\subsection{Fatou 引理、控制收敛与积分的绝对连续性}
\begin{theorem}[Fatou 引理]\label{thm:y0-fatou}
若 $f_n\ge0$ 可测，则
\[
 \int\liminf_nf_n\,\dd\mu\le\liminf_n\int f_n\,\dd\mu.
\]
若实值可积函数 $f_n$ 满足 $f_n\ge h$ a.e.，其中 $h$ 可积，则同样有 $\int\liminf_nf_n\le\liminf_n\int f_n$，左侧按有可积下界的扩展积分解释。若有可积上界，则得到上极限的反向不等式。
\end{theorem}
\begin{proof}
令 $g_N=\inf_{n\ge N}f_n$，则 $0\le g_N\uparrow\liminf_nf_n$，且 $\int g_N\le\inf_{n\ge N}\int f_n$。单调收敛给出结论。有可积下界时对 $f_n-h$ 应用非负版本，再加回 $\int h$；有上界时对 $-f_n$ 应用下界版本。
\end{proof}

\begin{theorem}[控制收敛定理]\label{thm:y0-dominated}
设 $f_n$ 与 $f$ 均可测，$f_n\to f$ a.e.，且存在非负可积函数 $g$ 使 $|f_n|\le g$ a.e.，则 $f$ 可积，$\int|f_n-f|\to0$，并且 $\int f_n\to\int f$。
\end{theorem}
\begin{proof}
有 $|f|\le g$ a.e.，故 $f$ 可积。函数 $h_n=|f_n-f|$ 趋零且 $0\le h_n\le2g$。对非负函数 $2g-h_n$ 使用 Fatou 引理，得
\[
 2\int g\le\liminf_n\left(2\int g-\int h_n\right)
            =2\int g-\limsup_n\int h_n.
\]
所以 $\int h_n\to0$，再用积分估计即可。
\end{proof}
这里把极限函数的可测性单独列出：在未完备的测度空间中，对可测零集上的值作任意修改，可能破坏可测性。若只给出 $f_n$ 几乎处处收敛，则应选取其可测极限代表，或先完备化测度空间，再使用上述结论。

\begin{proposition}[积分的可数可加性与绝对连续性]\label{prop:y0-integral-ac}
若 $f$ 可积，$\nu_f(E)=\int_Ef\,\dd\mu$ 为可数可加的实值或复值集合函数，且 $\mu(E)=0$ 蕴含 $\nu_f(E)=0$。更强地，对任意 $\varepsilon>0$，存在 $\delta>0$，使 $\mu(E)<\delta$ 时 $\int_E|f|\,\dd\mu<\varepsilon$。
\end{proposition}
\begin{proof}
对不交 $E_n$，部分并的示性函数逐点趋于全并的示性函数，乘以 $f$ 后由 $|f|$ 控制，控制收敛给出可数可加性。选 $M\ge1$ 使 $\int_{\{|f|>M\}}|f|<\varepsilon/2$，这由控制收敛得到。于是
\[
 \int_E|f|\le M\mu(E)+\int_{\{|f|>M\}}|f|.
\]
取 $\delta=\varepsilon/(2M)$ 即得结论，零测集的结论也直接成立。
\end{proof}

\subsection{乘积测度与 Fubini--Tonelli 定理}
\begin{lemma}[$\pi$--$\lambda$ 原理]\label{lem:y0-pi-lambda}
集合族 $\mathcal P$ 若对有限交封闭，称为 $\pi$-系。集合族 $\mathcal D$ 若含全集、对补集及可数不交并封闭，称为 $\lambda$-系。若 $\mathcal P\subseteq\mathcal D$，则 $\sigma(\mathcal P)\subseteq\mathcal D$。
\end{lemma}
\begin{proof}
令 $\mathcal L$ 为包含 $\mathcal P$ 的最小 $\lambda$-系。$\lambda$-系对两个嵌套成员的差封闭。固定 $A\in\mathcal P$，集合族 $\{B:A\cap B\in\mathcal L\}$ 为 $\lambda$-系并包含 $\mathcal P$，所以包含 $\mathcal L$。再固定 $B\in\mathcal L$，同样论证 $\{A:A\cap B\in\mathcal L\}$ 包含 $\mathcal L$。于是 $\mathcal L$ 对有限交封闭，结合补集可得有限并封闭，再将可数并分为不交差集，得它对可数并封闭。因此 $\mathcal L$ 是 $\sigma$-代数，必包含 $\sigma(\mathcal P)$，结论成立。
\end{proof}

\begin{theorem}[乘积测度与 Tonelli 定理]\label{thm:y0-tonelli}
设 $(S,\Sigma,\mu)$、$(T,\mathcal T,\nu)$ 均 $\sigma$-有限。令 $\Sigma\otimes\mathcal T$ 为由矩形 $A\times B$ 生成的 $\sigma$-代数。存在唯一测度 $\mu\times\nu$ 满足
\[
 (\mu\times\nu)(A\times B)=\mu(A)\nu(B).
\]
对非负乘积可测函数 $f$，截面可测，两个截面积分函数可测，且
\[
 \int_{S\times T}f\,\dd(\mu\times\nu)
 =\int_S\left(\int_Tf(s,t)\,\dd\nu(t)\right)\dd\mu(s)
 =\int_T\left(\int_Sf(s,t)\,\dd\mu(s)\right)\dd\nu(t),
\]
允许值为无穷，乘积中的 $0\cdot\infty$ 约定为零。
\end{theorem}
\begin{proof}
先设两测度有限。对集合 $E$，记 $E_s=\{t:(s,t)\in E\}$、$E^t=\{s:(s,t)\in E\}$。全部截面可测的集合构成包含矩形的 $\sigma$-代数，故乘积可测集均有可测截面。

考虑满足以下性质的集合 $E$：$s\mapsto\nu(E_s)$、$t\mapsto\mu(E^t)$ 可测，且它们的积分相等。矩形具有这些性质。有限测度保证取补集时可将截面测度及其积分从有限总量中相减；取可数不交并时由单调收敛得到性质。因此此族为包含矩形的 $\lambda$-系，由 $\pi$--$\lambda$ 原理，性质对所有乘积可测集成立。定义
\[
 \lambda(E)=\int_S\nu(E_s)\,\dd\mu(s).
\]
对不交集族，截面不交，测度可加性与单调收敛证明 $\lambda$ 为测度，并有所需矩形值。它也等于反向截面积分。

一般 $\sigma$-有限情形，将 $S$、$T$ 各分为可数个不交的有限测度集 $S_i$、$T_j$。在每个 $S_i\times T_j$ 上用上述有限测度构造，再对 $(i,j)$ 求和，得全空间测度。截面函数的可测性及两向积分相等由非负级数的单调收敛继承。对任意非负简单函数先按示性函数公式计算，再以递增简单函数逼近 $f$，三次使用单调收敛便得 Tonelli 公式。

唯一性在每个有限测度块上由 $\pi$--$\lambda$ 原理得到：两个有限测度相等的集合构成 $\lambda$-系并包含矩形。再对所有块求和，得全空间唯一性。
\end{proof}

\begin{theorem}[Fubini 定理]\label{thm:y0-fubini}
在上一设定下，若乘积可测函数 $f$ 满足 $\int|f|\,\dd(\mu\times\nu)<\infty$，则几乎所有截面可积，截面积分函数可积，且两个累次积分均等于 $\int f\,\dd(\mu\times\nu)$。绝对可积性等价于任意一个 $|f|$ 的累次积分有限。
\end{theorem}
\begin{proof}
对 $|f|$ 应用 Tonelli 定理，截面的绝对积分几乎处处有限，否则其积分不能有限。对实部与虚部的正、负部分分别使用 Tonelli，再作有限量的相减，得到交换次序的公式；由 $|\int f(s,t)\,\dd\nu(t)|\le\int|f(s,t)|\,\dd\nu(t)$ 得截面积分可积。最后的等价性正是 $|f|$ 的 Tonelli 公式。
\end{proof}

对完备化后的乘积测度，可用与原乘积可测函数几乎处处相同的代表元解释这些结论，截面结论仍是几乎处处的。不能只因两个带符号的累次积分各自存在，就跳过绝对可积性的检查。

\subsection{拓扑测度、Borel 集与 Baire 集}
\begin{definition}[Borel、Baire 与 Radon 测度]\label{def:y0-topological-measure}
拓扑空间 $X$ 的 Borel $\sigma$-代数 $\mathcal B(X)$ 由开集生成。现代 Baire $\sigma$-代数由所有连续实值函数的开集原像生成，通常包含于 $\mathcal B(X)$。在局部紧 Hausdorff 空间上，若 Borel 测度在紧集上有限，并满足
\[
 \mu(E)=\inf_{\substack{U\supseteq E\\U\text{ 开}}}\mu(U),\qquad
 \mu(E)=\sup_{\substack{K\subseteq E\\K\text{ 紧}}}\mu(K)
 \quad(E\in\mathcal B(X)),
\]
则本稿称其为 Radon 测度；这一定义明确要求外正则性与内正则性。
\end{definition}

原书的 Baire 集从紧 $G_\delta$ 集生成 $\sigma$-环，其中 $G_\delta$ 表示可数个开集的交；Borel 集也使用紧集生成的表述。这在一般非 $\sigma$-紧空间中与现代全空间 $\sigma$-代数表述须区别对待。在 $\R^n$ 及其闭子空间中，上述区别消失：这些空间可由可数个紧集覆盖，每个闭集为 $G_\delta$，而开集可写为可数个紧集的并，因此原书两种可测集族均等于通常的 Borel 集族。对连续函数的可测性，这些欧氏空间中也有 Baire 与 Borel 一致的结论。

例如，闭集 $F$ 可写为 $\bigcap_m\{x:d(x,F)<1/m\}$；开集 $U\subseteq\R^n$ 可写为紧集
\[
 K_m=\{x:|x|\le m,\ d(x,\R^n\setminus U)\ge1/m\}
\]
的并。更一般拓扑空间中的正则性和 Baire--Borel 延拓，须使用相应的测度表示定理，不能仅由可数可加性推出。

\subsection{Lebesgue 测度与两个逼近性质}
\begin{definition}[Lebesgue 外测度与可测集]\label{def:y0-lebesgue}
对 $E\subseteq\R^n$，令
\[
 m^*(E)=\inf\left\{\sum_j|Q_j|:
               E\subseteq\bigcup_jQ_j,\ Q_j\text{ 为有界开矩形}\right\},
\]
其中 $|Q_j|$ 是边长的乘积。若对每个 $A\subseteq\R^n$ 有
\[
 m^*(A)=m^*(A\cap E)+m^*(A\setminus E),
\]
称 $E$ 为 Lebesgue 可测集；$m^*$ 在这些集合上的限制记为 $m$。
\end{definition}

\begin{proposition}[外测度构造的基本性质]\label{prop:y0-lebesgue-construction}
上述可测集形成完备 $\sigma$-代数，$m$ 为 $\sigma$-有限测度，所有 Borel 集均可测，矩形的测度等于其体积。$m$ 平移不变，且对有限测度可测集 $E$ 有
\[
 m(E)=\inf_{U\supseteq E,\ U\text{ 开}}m(U)
     =\sup_{K\subseteq E,\ K\text{ 紧}}m(K).
\]
\end{proposition}
\begin{proof}
覆盖定义直接给出 $m^*(\varnothing)=0$、单调性及可数次可加性；后一性质通过对各集合选误差小于 $\varepsilon2^{-j}$ 的覆盖，再把覆盖合并得到。

分割等式对补集对称。若 $E,F$ 满足分割等式，对任意测试集先按 $E$、再按 $F$ 分割，得四项分解，从而得 $E\cup F$ 也满足等式。对不交可测集 $E_j$，重复分割给出
\[
 m^*(A)\ge\sum_{j=1}^Nm^*(A\cap E_j)
                   +m^*\left(A\setminus\bigcup_{j\ge1}E_j\right).
\]
令 $N\to\infty$，结合 $m^*(A\cap\bigcup_jE_j)\le\sum_jm^*(A\cap E_j)$，得全并满足分割等式；将一般并化为不交差集，证明 $\sigma$-代数性质。取 $A=\bigcup_jE_j$，同一估计与次可加性给出可数可加性。零外测度集及其子集自动满足分割等式，故测度完备。

坐标半空间可测：对任意矩形覆盖，沿对应坐标超平面切开各矩形，可将交叠处增宽任意小的宽度，使两边仍由开矩形覆盖且总体积增加任意小。于是 $m^*(A\cap H)+m^*(A\setminus H)\le m^*(A)$；反向由次可加性成立。因此矩形可测，开集为可数个有理矩形的并，故 Borel 集可测。

矩形的外测度不超过其体积，可通过稍增宽的开矩形覆盖得到。反向，对矩形内的任意紧闭矩形，用任意可数开矩形覆盖，紧性给出有限子覆盖；按有限个边界坐标细分，有限覆盖的体积和不小于被覆盖闭矩形的体积。令内矩形扩张至原矩形，得反向不等式。平移保持矩形体积并一一对应覆盖，故外测度及可测测度平移不变。可数个有界矩形覆盖全空间，故 $\sigma$-有限。

对有限测度 $E$，取体积总和小于 $m(E)+\varepsilon$ 的开矩形覆盖，并集 $U$ 开，且 $m(U\setminus E)<\varepsilon$，得外正则性。为得内正则性，先取足够大的闭立方体 $Q$，使 $m(E\setminus Q)<\varepsilon$；对 $Q\setminus E$ 外逼近为开集 $V$，使 $m(V\setminus(Q\setminus E))<\varepsilon$。则 $K=Q\setminus V$ 紧、$K\subseteq E$，且 $m(E\setminus K)<2\varepsilon$。
\end{proof}

在实线上，单调不减且右连续函数 $F$ 决定的区间值 $\mu_F((a,b])=F(b)-F(a)$ 称为 Lebesgue--Stieltjes 测度的起点；一般 $F$ 可用区间覆盖的外测度构造完成延拓。\emph{只有 $F(t)=t$ 时才是通常的 Lebesgue 测度}，例如阶跃函数决定点质量。这里不把一般 Stieltjes 测度统称为 Lebesgue 测度。

\begin{theorem}[紧支集连续函数的 $L^1$ 稠密性]\label{thm:y0-cc-density}
若 $G\subseteq\R^n$ 开，$f\in L^1(G)$，则对任意 $\varepsilon>0$，存在 $\varphi\in C_c(G)$ 使 $\int_G|f-\varphi|\,\dd x<\varepsilon$。
\end{theorem}
\begin{proof}
由命题 \ref{prop:y0-simple-l1}，先将 $f$ 逼近为可积简单函数，只需处理 $\mathbf1_E$，其中 $E\subseteq G$ 可测且有限测度。由正则性取紧集 $K\subseteq E$ 和开集 $U\supseteq E$，使 $m(E\setminus K)$、$m(U\setminus E)$ 都任意小。将 $U$ 换为 $U\cap G$，再在其中选含 $K$ 的相对紧开邻域 $V$。由推论 \ref{cor:y0-cutoff} 取 $0\le\varphi\le1$，在 $K$ 上为 $1$ 且支集包含于 $V$。于是
\[
 \int_G|\mathbf1_E-\varphi|
 \le m(E\setminus K)+m(U\setminus E).
\]
将有限个示性函数的逼近按系数组合，得到结论；有限个紧支集的并仍紧且包含于 $G$。
\end{proof}

\begin{theorem}[有限测度集的平移连续性]\label{thm:y0-translation-set}
若 $E\subseteq\R^n$ 可测且 $m(E)<\infty$，则 $m((E+h)\mathbin{\triangle}E)\to0$ 随 $h\to0$，其中 $\triangle$ 为对称差。
\end{theorem}
\begin{proof}
左侧等于 $\|\mathbf1_E(\cdot-h)-\mathbf1_E\|_{L^1}$。取 $\varphi\in C_c(\R^n)$ 使 $\|\mathbf1_E-\varphi\|_1<\varepsilon$。平移不变性给出
\[
 \|\mathbf1_E(\cdot-h)-\mathbf1_E\|_1
 \le2\varepsilon+\|\varphi(\cdot-h)-\varphi\|_1.
\]
当 $|h|\le1$ 时，后者的支集处于一个固定紧集内；紧支集连续函数在全空间一致连续，所以其平移差一致趋零，积分也趋零。最后令 $\varepsilon\downarrow0$。
\end{proof}

\begin{exercise}[逐点收敛与积分收敛]
在 $(0,1)$ 上令 $f_n=n\mathbf1_{(0,1/n)}$。证明它几乎处处趋于零，但积分不趋于零，并说明它不能由一个可积函数统一控制。
\end{exercise}
\begin{solution}
固定 $x>0$，充分大时 $1/n<x$，故 $f_n(x)=0$，而 $\int_0^1f_n=1$。若有可积 $g\ge|f_n|$ a.e.，控制收敛定理将给出积分趋零，矛盾。也可由积分的绝对连续性直接看出：$\int_{(0,1/n)}g\ge1$ 与这些集合测度趋零矛盾。
\end{solution}

\begin{exercise}[非绝对可积时累次积分可以不同]
在 $(0,1)^2$ 上取 $f(x,y)=(x^2-y^2)/(x^2+y^2)^2$。计算两个累次积分，说明为什么不违反 Fubini 定理。
\end{exercise}
\begin{solution}
由于 $\partial_y\bigl(y/(x^2+y^2)\bigr)=f$，对固定 $x>0$ 有 $\int_0^1f(x,y)\,\dd y=1/(1+x^2)$，再积分得 $\pi/4$。由于 $\partial_x\bigl(x/(x^2+y^2)\bigr)=-f$，反向累次积分为 $-\pi/4$。在第一象限极坐标下，$|f|=|\cos(2\theta)|/r^2$，面积元为 $r\,\dd r\,\dd\theta$，原点附近径向积分 $\int_0^\delta\dd r/r$ 发散，故 $f$ 不绝对可积，Fubini 的假设不成立。
\end{solution}


\Needspace{8\baselineskip}
\section{线性空间}\label{sec:yosida-00-04}
\noindent\textit{原书位置：中译本与英文版均为印刷第 20--22 页。}
% !TeX encoding = UTF-8
本节只建立代数结构，不给线性空间预先指定拓扑。后续的范数、半范数、闭子空间及连续算子，都是在这些代数对象上添加分析条件。除特别说明外，标量域 $\mathbb K$ 为 $\R$ 或 $\mathbb C$。

\subsection{线性空间、子空间与生成空间}
\begin{definition}[线性空间]\label{def:y0-vector-space}
域 $\mathbb K$ 上的线性空间是集合 $V$，配有加法与数乘，满足：加法使 $V$ 成为交换群，零元记为 $0$；且对 $x,y\in V$、$\alpha,\beta\in\mathbb K$ 有
\[
 \alpha(x+y)=\alpha x+\alpha y,\qquad
 (\alpha+\beta)x=\alpha x+\beta x,\qquad
 (\alpha\beta)x=\alpha(\beta x),\qquad 1x=x.
\]
其元素称为向量。
\end{definition}

由分配律与加法消去律可得 $0x=0$、$\alpha0=0$、$(-1)x=-x$；若 $\alpha x=0$ 且 $\alpha\ne0$，乘以 $\alpha^{-1}$ 得 $x=0$。标量零与向量零通常共用记号，语境决定其所属空间。

\begin{definition}[子空间与线性张成]\label{def:y0-subspace}
非空子集 $M\subseteq V$ 若对加法与数乘封闭，称为线性子空间。等价地，对任意 $x,y\in M$ 及标量 $\alpha,\beta$，有 $\alpha x+\beta y\in M$。

集合 $A\subseteq V$ 的线性张成 $\operatorname{span}A$ 是 $A$ 中有限个向量的全部线性组合组成的集合，约定 $\operatorname{span}\varnothing=\{0\}$。
\end{definition}

\begin{proposition}[张成与子空间的运算]\label{prop:y0-span}
$\operatorname{span}A$ 是包含 $A$ 的最小子空间。任意族子空间的交仍为子空间，空族的交为 $V$。若 $M,N$ 为子空间，则 $M+N=\{m+n:m\in M,n\in N\}$ 是包含 $M\cup N$ 的最小子空间。
\end{proposition}
\begin{proof}
有限线性组合的和或数乘仍为有限线性组合，故张成是子空间；任何包含 $A$ 的子空间必须包含这些组合，得最小性。交中的元素在各成员中都对线性组合封闭。对 $M+N$，线性组合可分别整理到 $M$ 和 $N$ 中，故它是子空间；包含两者的任一子空间都必须包含它们的和。
\end{proof}

子空间的并通常不是子空间。例如 $\R^2$ 中两条不同的过原点直线，其并不含从两直线各取一个非零向量所得的和。若一族子空间按包含关系成链，则其并为子空间，因为任意两个向量共同属于某个较大的成员。

\begin{example}[后续使用的线性空间]
$\mathbb K^n$、全体数列 $\mathbb K^{\N}$、任意集合上的函数空间 $\mathbb K^S$ 都以逐坐标运算成为线性空间。$C(X;\mathbb K)$、$C_c(X;\mathbb K)$ 是函数空间的子空间；$L^p$ 则把几乎处处相等的可测函数识别为同一个向量。后一识别是一个商构造，它保证范数为零的对象恰为零向量。
\end{example}

\subsection{线性无关、基与维数}
\begin{definition}[线性无关与 Hamel 基]\label{def:y0-linear-independence}
有限向量族 $v_1,\ldots,v_n$ 线性无关，指 $\sum_j\alpha_jv_j=0$ 蕴含全部 $\alpha_j=0$。任意子集 $B\subseteq V$ 线性无关，指其中每个有限的互异向量族都线性无关。若 $B$ 同时张成 $V$，称为 Hamel 基。

若 $V$ 有 $n$ 个向量组成的基，称 $V$ 为 $n$ 维，记 $\dim V=n$；若没有有限基，称为无限维。零空间的空基给出维数 $0$。
\end{definition}

\begin{lemma}[有限交换引理]\label{lem:y0-exchange}
若 $u_1,\ldots,u_m$ 线性无关，而 $v_1,\ldots,v_n$ 张成 $V$，则 $m\le n$；且可用 $u_1,\ldots,u_m$ 替换生成族中的 $m$ 个向量而仍张成 $V$。
\end{lemma}
\begin{proof}
把 $u_1$ 表示为 $v_j$ 的线性组合。因为 $u_1\ne0$，有一个系数非零，可解出对应的 $v_j$，从而用 $u_1$ 替换它，仍保持张成。若已将 $u_1,\ldots,u_k$ 换入，把 $u_{k+1}$ 表示为当前生成族的线性组合；由于 $u_{k+1}$ 不属于前 $k$ 个 $u$ 的张成，至少有一个尚未替换的 $v_j$ 的系数非零，再解出它并替换。若 $m>n$，到第 $n$ 步时已无剩余 $v_j$，$u_{n+1}$ 将属于前 $n$ 个 $u$ 的张成，矛盾。
\end{proof}

\begin{corollary}[有限维空间的基定理]\label{cor:y0-finite-basis}
有限维空间中，任意两个基的元素个数相同，任意线性无关族可扩充为基，任意有限生成族可删减为基。在 $n$ 维空间中，任意 $n$ 个线性无关向量都是基，任意 $n$ 个生成向量也都是基。
\end{corollary}
\begin{proof}
对两个基分别用交换引理，得其元素个数互相不超过，从而相同。从有限生成族中逐一删去可由其余向量表示者，最终得到基；交换引理先将无关族换入有限生成族，再只删去多余的其余向量，即得到包含原无关族的基。若 $n$ 个无关向量未张成，可再加入一个不在其张成中的向量，与交换引理矛盾。若 $n$ 个生成向量相关，可删去至少一个，导致由少于 $n$ 个向量张成空间，也与交换引理矛盾。
\end{proof}

任意维空间中基的存在性已在定理 \ref{thm:y0-hamel-basis} 中用 Zorn 引理证明。这里有限维的结论只用有限交换，不需要一般的选择公理。具有基时，每个向量的有限坐标表达唯一，因两个表达相减给出有限线性关系。

\subsection{线性算子、核、值域与延拓}
\begin{definition}[线性算子与线性泛函]\label{def:y0-linear-operator}
设 $X,Y$ 为同一标量域上的线性空间，$D(T)\subseteq X$ 为子空间。映射 $T:D(T)\to Y$ 若满足
\[
 T(\alpha x+\beta y)=\alpha Tx+\beta Ty,
\]
称为线性算子。定义域、值域与核分别记为
\[
 D(T),\qquad R(T)=T[D(T)],\qquad N(T)=\ker T=\{x\in D(T):Tx=0\}.
\]
若 $Y=\mathbb K$，称 $T$ 为线性泛函。双射的线性算子称为线性同构。
\end{definition}

\begin{proposition}[核、值域与逆算子]\label{prop:y0-operator}
核及值域均为子空间。$T$ 单射当且仅当 $N(T)=\{0\}$；此时 $T^{-1}:R(T)\to D(T)$ 线性。映射 $T$ 若要在整个 $Y$ 上具有逆，还必须满足 $R(T)=Y$。
\end{proposition}
\begin{proof}
线性性给出 $T0=0$，以及核和值域对线性组合封闭。因为 $Tx=Ty$ 当且仅当 $T(x-y)=0$，单射性等价于零核。单射时，对 $u=Tx$、$v=Ty$，有 $\alpha u+\beta v=T(\alpha x+\beta y)$，故 $T^{-1}(\alpha u+\beta v)=\alpha x+\beta y$，证明逆算子线性。整个 $Y$ 上的逆需对每个 $y$ 有原像，即满射。
\end{proof}

两个算子相等，要求定义域相同且逐点取值相同。若 $D(T_1)\subseteq D(T_2)$ 且 $T_2|_{D(T_1)}=T_1$，称 $T_2$ 为 $T_1$ 的延拓，记 $T_1\subseteq T_2$；这也是将算子看作图时的包含关系。线性泛函的配对在本稿中对向量线性，例如 $\langle T,x\rangle=T(x)$；若使用内积实现泛函，则须遵守第一变量线性的既定内积约定。

\begin{proposition}[由基上的取值定义算子]\label{prop:y0-basis-map}
设 $B$ 为 $X$ 的 Hamel 基。任意映射 $a:B\to Y$ 唯一延拓为线性算子 $T:X\to Y$，由
\[
 T\left(\sum_{b\in F}\alpha_bb\right)
          =\sum_{b\in F}\alpha_ba(b)
 \qquad(F\subseteq B\text{ 有限})
\]
定义。
\end{proposition}
\begin{proof}
基的有限坐标表达唯一，故定义不依赖表达；将两向量的表达统一到一个有限基向量集合上，即可验证线性性。任何延拓都须满足同一公式，故唯一。
\end{proof}

这是代数构造；没有连续性结论。后续分析中必须另行检查定义的算子是否有界或连续。

\subsection{商空间与第一同构定理}
\begin{theorem}[商线性空间]\label{thm:y0-quotient}
设 $M\subseteq X$ 为子空间。关系 $x\sim y$ 当且仅当 $x-y\in M$ 是等价关系，其等价类为 $[x]=x+M$。在集合 $X/M=\{x+M:x\in X\}$ 上定义
\[
 [x]+[y]=[x+y],\qquad\alpha[x]=[\alpha x].
\]
这些运算良定，使 $X/M$ 成为线性空间。商映射 $q:X\to X/M$、$q(x)=[x]$ 为满射线性算子，且 $\ker q=M$。
\end{theorem}
\begin{proof}
$x-x=0\in M$ 给出自反性；$x-y\in M$ 蕴含其负向量 $y-x\in M$，给出对称性；两个差相加给出传递性。若 $x'=x+m$、$y'=y+n$，则 $(x'+y')-(x+y)=m+n\in M$，且 $\alpha x'-\alpha x=\alpha m\in M$，所以运算不依赖代表元。线性空间公理由原空间逐条传递，零元为 $M=[0]$，负元为 $[-x]$。商映射的线性性和满射性直接由定义成立；$[x]=[0]$ 当且仅当 $x\in M$，故核为 $M$。
\end{proof}

\begin{theorem}[商空间的因子分解与第一同构定理]\label{thm:y0-first-isomorphism}
设 $T:X\to Y$ 线性，$M\subseteq\ker T$。则存在唯一线性算子 $\widetilde T:X/M\to Y$ 使 $T=\widetilde T\circ q$，其公式为 $\widetilde T([x])=Tx$。特别地，
\[
 X/\ker T\ \cong\ R(T).
\]
对定义在子空间上的算子，相应结论为 $D(T)/N(T)\cong R(T)$。
\end{theorem}
\begin{proof}
若 $[x]=[y]$，则 $x-y\in M\subseteq\ker T$，所以 $Tx=Ty$，定义良定。线性性由商运算继承，唯一性由 $q$ 满射得到。取 $M=\ker T$ 并将目标限制为 $R(T)$，所得算子满射，且 $\widetilde T([x])=0$ 蕴含 $x\in\ker T$、$[x]=0$，所以为同构。
\end{proof}

\begin{corollary}[有限维的秩--零度公式]\label{cor:y0-rank-nullity}
若 $X$ 有限维且 $T:X\to Y$ 线性，则 $\dim X=\dim\ker T+\dim R(T)$。
\end{corollary}
\begin{proof}
取核的一组基 $e_1,\ldots,e_k$ 并扩充为 $X$ 的基 $e_1,\ldots,e_k,u_1,\ldots,u_r$。向量 $Tu_1,\ldots,Tu_r$ 张成值域；若它们有线性关系，则对应的 $u_j$ 组合属于核，再由原基的线性无关性得系数全零。因此它们为值域的基，维数公式成立。
\end{proof}

\begin{example}[商空间消除积分中的代表元歧义]
在可积函数构成的线性空间中，令 $M$ 为几乎处处为零的函数组成的子空间。商空间把 $f$ 与 $g$ 识别为同一元素，恰当且仅当 $f=g$ a.e.。积分距离 $\int|f-g|$ 对代表元不敏感，并且在商空间中只有零元素的距离为零，这就是 $L^1$ 空间的代数起点。抽象商空间无需子空间闭；要得到 Hausdorff 的拓扑商或真正的商范数，则需后续的闭性条件。
\end{example}

\begin{exercise}[代数补空间与投影]
假设选择公理。证明每个子空间 $M\subseteq X$ 都存在子空间 $N$，使 $X=M\oplus N$，即每个 $x$ 唯一写成 $m+n$。由此构造投影 $P:X\to M$，并说明它未必连续。
\end{exercise}
\begin{solution}
取 $M$ 的 Hamel 基并用定理 \ref{thm:y0-hamel-basis} 扩充为 $X$ 的基。令 $N$ 为新增基向量的张成。基的唯一有限表达给出 $X=M+N$ 及 $M\cap N=\{0\}$。定义 $P(m+n)=m$，则 $P$ 线性、$P^2=P$、$R(P)=M$、$\ker P=N$。例如在 $\ell^2$ 中，取包含全部单位向量 $e_n$ 的 Hamel 基，按命题 \ref{prop:y0-basis-map} 定义泛函 $f(e_n)=n$，其余基向量取值为零。它在单位向量上无界，因而不连续；对应的投影 $P(x)=f(x)e_1$ 满足 $P^2=P$，同样不连续。
\end{solution}


% !TeX encoding = UTF-8
\section*{本章参考文献与阅读衔接}
\phantomsection
\addcontentsline{toc}{section}{本章参考文献与阅读衔接}

本章主线依据 \cite{yosida-zh,yosida-en}。原书第零章末的参考条目已核对：中译本与英文版均位于印刷第 22 页，即两份本地 PDF 的第 37 页。它按主题列出拓扑空间的三组参考作者和测度空间的两组参考作者，没有另列集合论与线性空间书目。原书的作者后缀“[1]”指各作者在全书参考书目中的第一条，\emph{不是本稿的参考文献编号}。对应关系如下。

\begin{enumerate}
 \item \textbf{P. Alexandroff--H. Hopf [1]}：\emph{Topologie}, Vol.~I, Springer, 1935，见 \cite{alexandroff-hopf}。英文原书总书目印刷第 469 页（PDF 第 484 页）使用拼写 Alexandrov；章末及中译本使用 Alexandroff，均指同一作者。
 \item \textbf{N. Bourbaki [1]}：\emph{Topologie G\'en\'erale}，Hermann，原书记录年份 1940--1942，见 \cite{bourbaki-topology}。核对位置为英文原书总书目印刷第 470 页（PDF 第 485 页）。这是分册出版的书目记录；不可把现代合订本的重印年份当作原书所引年份。
 \item \textbf{J. L. Kelley [1]}：\emph{General Topology}, D.~Van Nostrand, 1955，见 \cite{kelley-topology}。核对位置为英文原书总书目印刷第 475 页（PDF 第 490 页）。Springer 的 GTM 27 版为后来的再版，书目中另作说明。
 \item \textbf{P. R. Halmos [1]}：\emph{Measure Theory}, D.~Van Nostrand, 1950，见 \cite{halmos-measure}。核对位置为英文原书总书目印刷第 473 页（PDF 第 488 页）。
 \item \textbf{S. Saks [1]}：\emph{Theory of the Integral}, Warszawa, 1937，见 \cite{saks-integral}。核对位置为英文原书总书目印刷第 481 页（PDF 第 496 页）。
\end{enumerate}

以上五条均已并入书后的统一参考文献表，正文不再维护另一套编号。这里记录的是吉田原书的参考指引；本次实际核对了原书章末与总书目的对应条目，未把这五部参考书的全文标为已经逐页查阅。

阅读衔接上，第 \ref{sec:yosida-00-01} 节的 Zorn 引理用于代数基及延拓构造；第 \ref{sec:yosida-00-02} 节的紧性、类别定理与连续函数逼近分别服务于弱拓扑、Baire 方法和表示定理；第 \ref{sec:yosida-00-03} 节的积分工具为 $L^p$ 空间、卷积及积分算子准备基础；第 \ref{sec:yosida-00-04} 节的核、值域和商空间则在后续算子论中反复出现。第一章从这些基础出发，以半范数建立拓扑向量空间。

